% GENERATED FILE. Edit canonical lessons, then run npm run generate:books.
% canonical-source-hash: fnv1a64:10a8fc8bfe4e0a06
% canonical-lessons: TA-C342-punaramai, TA-C342-pori-engrave, TA-C342-pantaiya, TA-C342-yukam, TA-C342-tolliyal

\chapter{Old Things Made New}
\label{ch:ta-old-things-made-new}
\hlchaptermodality{342}

\begin{chapteropening}
\textbf{By the end of this chapter:} \emph{I can talk about restoring and studying ancient places.}

Complete the last lesson of chapter 342: I can talk about restoring and studying ancient places.
\end{chapteropening}

\section[\texorpdfstring{\ta{புனரமை} (puṉaramai)}{புனரமை (puṉaramai)}]{\ta{புனரமை} (puṉaramai) — to restore; to renovate}

\label{lesson:TA-C342-punaramai}



\begin{quote}
Before the new one: say the Tamil for to dig out, then the Tamil for to come upon.
\end{quote}

\subsection*{You'll want to know: \ta{புனரமை}}
\textbf{\ta{புனரமை}} — \emph{puṉaramai} — \textquotedblleft{}to restore; to renovate\textquotedblright{}.

\begin{quote}
\ta{அரசு} \ta{அந்தப்} \ta{பழைய} \ta{கோட்டையைப்} \ta{புனரமைத்தது}.
\end{quote}

\begin{quote}
\emph{aracu antap paḻaiya kōṭṭaiyaip puṉaramaittatu.}
\end{quote}

\begin{quote}
The government restored that old fort.
\end{quote}

Past (I …): \ta{புனரமைத்தேன்}, \emph{puṉaramaittēṉ}.

\ta{புனர்}, again + \ta{அமை}, to build.

\begin{cousinweb}
\textbf{\ta{புனரமை}} (\emph{puṉaramai}) joins Sanskrit \emph{punar}, \textquotedblleft{}back, again\textquotedblright{} (CDIAL 8273), with native \ta{அமை}, \textquotedblleft{}to be settled, built\textquotedblright{} (DEDR 161).
\end{cousinweb}

\begin{grammarlens}[title={What you've built}]
One more word for telling what happened.
\end{grammarlens}

\subsection*{Your turn}
\begin{itemize}
\raggedright
  \item \emph{Say it:} \emph{puṉaramai}
  \item \emph{Say it:} \emph{puṉaramai}, once more
  \item \emph{Say it:} \emph{aracu antap paḻaiya kōṭṭaiyaip puṉaramaittatu.} — The government restored that old fort.
  \item \emph{Recall:} say the Tamil for to dig out, then the Tamil for to come upon, then say \emph{puṉaramai} again
\end{itemize}

\begin{tcolorbox}[breakable,colback=teal!4,colframe=teal!35!black,title={Before you move on}]
What does \ta{புனரமை} mean? (\textquotedblleft{}To restore; to renovate\textquotedblright{}.) Say it once more.
\end{tcolorbox}
\section[\texorpdfstring{\ta{பொறி} (poṟi)}{பொறி (poṟi)}]{\ta{பொறி} (poṟi) — to engrave; to inscribe}

\label{lesson:TA-C342-pori-engrave}



\begin{quote}
Before the new one: say the Tamil for to come upon, then the Tamil for to restore.
\end{quote}

\subsection*{You'll want to know: \ta{பொறி}}
\textbf{\ta{பொறி}} — \emph{poṟi} — \textquotedblleft{}to engrave; to inscribe\textquotedblright{}.

\begin{quote}
\ta{அவர்கள்} \ta{கல்லில்} \ta{மன்னரின்} \ta{பெயரைப்} \ta{பொறித்தார்கள்}.
\end{quote}

\begin{quote}
\emph{avarkaḷ kallil maṉṉariṉ peyaraip poṟittārkaḷ.}
\end{quote}

\begin{quote}
They engraved the king's name on the stone.
\end{quote}

Past (I …): \ta{பொறித்தேன்}, \emph{poṟittēṉ}.

The place takes -\ta{இல்}: \ta{கல்லில்} \ta{பொறி}.

\begin{cousinweb}
\textbf{\ta{பொறி}} (\emph{poṟi}) is native Dravidian (DEDR 4564): to impress, stamp, inscribe.
\end{cousinweb}

\begin{grammarlens}[title={What you've built}]
One more word for telling what happened.
\end{grammarlens}

\subsection*{Your turn}
\begin{itemize}
\raggedright
  \item \emph{Say it:} \emph{poṟi}
  \item \emph{Say it:} \emph{poṟi}, once more
  \item \emph{Say it:} \emph{avarkaḷ kallil maṉṉariṉ peyaraip poṟittārkaḷ.} — They engraved the king's name on the stone.
  \item \emph{Recall:} say the Tamil for to come upon, then the Tamil for to restore, then say \emph{poṟi} again
\end{itemize}

\begin{tcolorbox}[breakable,colback=teal!4,colframe=teal!35!black,title={Before you move on}]
What does \ta{பொறி} mean? (\textquotedblleft{}To engrave; to inscribe\textquotedblright{}.) Say it once more.
\end{tcolorbox}
\section[\texorpdfstring{\ta{பண்டைய} (paṇṭaiya)}{பண்டைய (paṇṭaiya)}]{\ta{பண்டைய} (paṇṭaiya) — ancient}

\label{lesson:TA-C342-pantaiya}



\begin{quote}
Before the new one: say the Tamil for to restore, then the Tamil for to engrave.
\end{quote}

\subsection*{You'll want to know: \ta{பண்டைய}}
\textbf{\ta{பண்டைய}} — \emph{paṇṭaiya} — \textquotedblleft{}ancient\textquotedblright{}.

\begin{quote}
\ta{இது} \ta{ஒரு} \ta{பண்டைய} \ta{நகரம்}.
\end{quote}

\begin{quote}
\emph{itu oru paṇṭaiya nakaram.}
\end{quote}

\begin{quote}
This is an ancient city.
\end{quote}

Only before a noun; from \ta{பண்டு}, former times.

\begin{cousinweb}
\textbf{\ta{பண்டைய}} (\emph{paṇṭaiya}) is native Dravidian: from \ta{பண்டை}, \ta{பண்டு}, \textquotedblleft{}antiquity, former time\textquotedblright{} (DEDR 3999).
\end{cousinweb}

\begin{grammarlens}[title={What you've built}]
One more word for telling what happened.
\end{grammarlens}

\subsection*{Your turn}
\begin{itemize}
\raggedright
  \item \emph{Say it:} \emph{paṇṭaiya}
  \item \emph{Say it:} \emph{paṇṭaiya}, once more
  \item \emph{Say it:} \emph{itu oru paṇṭaiya nakaram.} — This is an ancient city.
  \item \emph{Recall:} say the Tamil for to restore, then the Tamil for to engrave, then say \emph{paṇṭaiya} again
\end{itemize}

\begin{tcolorbox}[breakable,colback=teal!4,colframe=teal!35!black,title={Before you move on}]
What does \ta{பண்டைய} mean? (\textquotedblleft{}Ancient\textquotedblright{}.) Say it once more.
\end{tcolorbox}
\section[\texorpdfstring{\ta{யுகம்} (yukam)}{யுகம் (yukam)}]{\ta{யுகம்} (yukam) — an age; an era}

\label{lesson:TA-C342-yukam}



\begin{quote}
Before the new one: say the Tamil for to engrave, then the Tamil for ancient.
\end{quote}

\subsection*{You'll want to know: \ta{யுகம்}}
\textbf{\ta{யுகம்}} — \emph{yukam} — \textquotedblleft{}an age; an era\textquotedblright{}.

\begin{quote}
\ta{இது} \ta{கணினி} \ta{யுகம்}.
\end{quote}

\begin{quote}
\emph{itu kaṇiṉi yukam.}
\end{quote}

\begin{quote}
This is the computer age.
\end{quote}

Also one of the long ages of Indian tradition.

\begin{cousinweb}
\textbf{\ta{யுகம்}} (\emph{yukam}) is borrowed from Sanskrit \emph{yuga}, \textquotedblleft{}yoke, a pair\textquotedblright{} (CDIAL 10482), which also came to mean an age of the world.
\end{cousinweb}

\begin{grammarlens}[title={What you've built}]
One more word for telling what happened.
\end{grammarlens}

\subsection*{Your turn}
\begin{itemize}
\raggedright
  \item \emph{Say it:} \emph{yukam}
  \item \emph{Say it:} \emph{yukam}, once more
  \item \emph{Say it:} \emph{itu kaṇiṉi yukam.} — This is the computer age.
  \item \emph{Recall:} say the Tamil for to engrave, then the Tamil for ancient, then say \emph{yukam} again
\end{itemize}

\begin{tcolorbox}[breakable,colback=teal!4,colframe=teal!35!black,title={Before you move on}]
What does \ta{யுகம்} mean? (\textquotedblleft{}An age; an era\textquotedblright{}.) Say it once more.
\end{tcolorbox}
\section[\texorpdfstring{\ta{தொல்லியல்} (tolliyal)}{தொல்லியல் (tolliyal)}]{\ta{தொல்லியல்} (tolliyal) — archaeology}

\label{lesson:TA-C342-tolliyal}



\begin{quote}
Before the new one: say the Tamil for ancient, then the Tamil for an age.
\end{quote}

\subsection*{You'll want to know: \ta{தொல்லியல்}}
\textbf{\ta{தொல்லியல்}} — \emph{tolliyal} — \textquotedblleft{}archaeology\textquotedblright{}.

\begin{quote}
\ta{அவள்} \ta{பல்கலைக்கழகத்தில்} \ta{தொல்லியல்} \ta{படிக்கிறாள்}.
\end{quote}

\begin{quote}
\emph{avaḷ palkalaikkaḻakattil tolliyal paṭikkiṟāḷ.}
\end{quote}

\begin{quote}
She studies archaeology at university.
\end{quote}

\ta{தொல்}, ancient + \ta{இயல்}, a field of study.

\begin{cousinweb}
\textbf{\ta{தொல்லியல்}} (\emph{tolliyal}) is a Tamil compound built from \ta{தொல்}, \textquotedblleft{}old, ancient\textquotedblright{} (DEDR 3516), and \ta{இயல்}, \textquotedblleft{}nature, a branch of learning\textquotedblright{}.
\end{cousinweb}

\begin{grammarlens}[title={What you've built}]
One more word for telling what happened.
\end{grammarlens}

\subsection*{Your turn}
\begin{itemize}
\raggedright
  \item \emph{Say it:} \emph{tolliyal}
  \item \emph{Say it:} \emph{tolliyal}, once more
  \item \emph{Say it:} \emph{avaḷ palkalaikkaḻakattil tolliyal paṭikkiṟāḷ.} — She studies archaeology at university.
  \item \emph{Recall:} say the Tamil for ancient, then the Tamil for an age, then say \emph{tolliyal} again
\end{itemize}

\begin{tcolorbox}[breakable,colback=teal!4,colframe=teal!35!black,title={Before you move on}]
What does \ta{தொல்லியல்} mean? (\textquotedblleft{}Archaeology\textquotedblright{}.) Say it once more.
\end{tcolorbox}
